%%%%%%%%%%%%%%%%%%%%%%%%%%%%%%%%%%%%%%%%%%%%%%%%%%%%%%%%%%%%
%% CHAPTER 6 -- Pretraining and the Birth of GPT
%%%%%%%%%%%%%%%%%%%%%%%%%%%%%%%%%%%%%%%%%%%%%%%%%%%%%%%%%%%%

\chapter{Pretraining and the Birth of GPT}
\label{ntg:ch6}

\section*{Setting the Scene}
\addcontentsline{toc}{section}{Setting the Scene}

\subsection*{The problem}

Picture an engineer in 2018, holding a fresh implementation of
the Transformer architecture from the previous chapter and a
hard drive full of crawled web text. The model runs. The text
is there. The question is what to do with the text in order to
get a model that is, in some useful sense, fluent.

The classical answer was to pick a narrow task --- machine
translation, sentiment classification, summarisation --- and
train the model on labelled examples of that task. The trouble
with this answer was that labelled examples are expensive (a
human has to produce them), that any single task is a thin
slice of what language can be used for, and that a model
trained only on translation has no obvious reason to be good at
anything else. What you would want, instead, is a single
training procedure that produced a single model that was good
at many things --- ideally, things you had not specifically
asked it to be good at.

The procedure that emerged is now called \emph{pretraining}. It
is the procedure of Chapter~1 (predict the next token) and
Chapter~2 (gradient descent) and Chapter~5 (Transformer),
applied at a scale almost no one anticipated. The result is a
single model that is, in some broad sense, capable of being
talked to about almost any subject the training data covered.
The same model, after the additional steps of Chapter~9, is
the assistant on the other side of every chat interface you
have used in the last three years.

\subsection*{How we got here}

The connection between language modelling and information
theory is older than the field. Claude Shannon, in a 1951
paper called \emph{Prediction and Entropy of Printed English},
estimated the predictability of the language by playing a
guessing game with human subjects. He showed them passages of
text with the next letter withheld, and recorded how often they
guessed correctly. The performance of the guessers gave him a
lower bound on how surprising the language is, and therefore on
how well any language model could ever do. He treated the human
guessers as the model. The connection between predicting the
next symbol and modelling the language has been with us ever
since.

What was missing, for half a century after Shannon, was a
sufficiently expressive learner and a sufficiently large
corpus. Both arrived together in the 2010s. Bengio's 2003 paper
demonstrated the principle of a neural language model. The
2013 word2vec work demonstrated that learned representations
of words could capture meaning in their geometry. The 2017
Transformer demonstrated that the attention mechanism could
support models large enough to use billions of parameters
productively. By 2018 the pieces were on the table.

The first GPT model --- the name stands, prosaically, for
Generative Pre-trained Transformer --- was published by Alec
Radford and his collaborators at OpenAI in June 2018. The
architecture was a Transformer decoder. The training procedure
was pretraining on a corpus called BookCorpus (about seven
thousand unpublished books). The model had a hundred and ten
million parameters. By modern standards it was tiny. The point
of the paper was that, once the model had been pretrained on
the books, fine-tuning it on a small amount of labelled data
for a downstream task gave better results than training a
specialist model from scratch on the same labelled data.
Pretraining had become the recommended starting point.

GPT-2, in February 2019, was the same idea at fifteen times the
parameter count and several times the data. Its release was
notable for what OpenAI did not do: they declined, for several
months, to release the largest version of the model, citing
concerns about possible misuse. The decision was widely
discussed and is, in retrospect, one of the early markers of the
field's growing sense that the artefacts it was building had
implications beyond the academic literature. They eventually
released the full model. Critics observed that the misuse
predictions had largely not materialised. The episode set a
precedent for a long series of similar decisions to come.

GPT-3, in May 2020, was the same idea at a hundred and fifteen
times the parameter count of GPT-2 --- one hundred and seventy-
five billion parameters --- trained on hundreds of billions of
tokens. The paper, by Tom Brown and a long list of coauthors,
was the one that made the rest of the field sit up. It
documented something that had not been the explicit aim of the
training: GPT-3, simply by being shown a few examples of a task
in its prompt, could perform that task without any further
training. It could do arithmetic. It could translate between
languages it had not been specifically trained to translate. It
could solve simple word problems. The effect was not perfect,
and it was much better at large scales than at small ones, but
it was robust enough to be reproduced and dramatic enough to
shift the field's expectations of what a pretrained model
could do. The authors called the phenomenon \emph{in-context
learning}. Two years later, the same model would be wrapped in
the alignment work of Chapter~9 to produce ChatGPT.

ChatGPT was released on 30 November 2022 as a research
preview. It crossed a hundred million users in two months ---
the fastest consumer product adoption of any kind in recorded
history. The release was, in a sense, the moment a
twenty-five-year arc of academic work suddenly became visible
to the public. GPT-4 followed in March 2023. Claude 1, 2, and
3 from Anthropic, Gemini from Google, Llama 1, 2, and 3 from
Meta, and DeepSeek and Qwen from Chinese labs filled out the
landscape over the following two years. The architectural
ideas in all of them are, with very few exceptions, the
Transformer of Chapter~5 trained by the procedure described in
this chapter.

\subsection*{Where we stand}

A modern pretraining run has the following character. The data
is a curated corpus on the order of trillions of tokens,
assembled from the public web, books, code repositories,
academic papers, and a varying mix of licensed proprietary
sources. The data has been filtered for quality (a great deal
of the web is repetitive, machine-generated, or simply low
information) and deduplicated (many documents on the web are
near-copies of one another). The data has been mixed in
proportions that the model's developers believe will produce
the desired balance of capabilities --- more code if the model
is meant to be good at programming, more academic text if the
model is meant to be good at scientific reasoning, more
conversational data if the model is meant to be good at
dialogue.

The model is a Transformer of the kind described in Chapter~5,
sized according to the Chinchilla scaling law of Chapter~2 to
match the available compute and data. It is trained, for
several weeks to several months, on a cluster of thousands of
graphics processors. The training procedure is the gradient
descent of Chapter~2 with the cross-entropy loss of Chapter~1.
At the end of the training, the model is a \emph{base model}:
fluent at completing text in the style of its training corpus,
capable of in-context learning, but not yet capable of behaving
as an assistant.

The transformation from a base model into an assistant is the
subject of Chapter~9. The base model is the raw material; the
assistant is the manufactured product. The two are different
artefacts, with different behaviours, and the language we use
to discuss them should reflect the difference.

\subsection*{What goes wrong}

Pretraining on internet-scale text has costs that the loss
function does not see.

The training data contains copyrighted material, which the
model can be induced to reproduce verbatim under the right
prompt. The training data contains personal information, which
the same is true of. The training data contains the prejudices
of its authors, which the model will reproduce fluently. The
training data contains misinformation, which the model will
sometimes reproduce confidently. The model has no native means
of distinguishing reliable from unreliable sources within its
training corpus, because the training procedure does not
include any signal about reliability --- only about how often a
given continuation appears.

The training corpus has a cut-off date. After that date, the
model knows nothing. The world continues to change. A model
trained on data up to October 2023 that is asked, in 2025,
about events after that date will produce a confident answer
that is, with respect to those events, fictional. The model has
no native means of telling the user that the question lies past
its cut-off. We will see in Chapter~10 how retrieval can fill
the gap, with caveats.

The phenomenon of in-context learning, while striking, is
poorly understood. We do not have, at this writing, a settled
account of why showing a model a handful of examples in its
prompt can change its behaviour as much as it does. We have a
few candidate explanations and a continuing argument among them.
The practical implication is that designing prompts that
reliably produce good in-context behaviour is more art than
science, and prompts that work well in one context can fail in
seemingly identical ones.

\section{The idea, in plain language}
\addcontentsline{toc}{section}{The idea, in plain language}

\subsection*{What ``pretraining'' actually means}

The word \emph{pretraining} is a tell. It implies that there is
a later stage, sometimes called fine-tuning or alignment, that
follows. Both implications are correct. The flow goes:

\begin{enumerate}
  \item \textbf{Pretraining.} A Transformer of the appropriate
        size is trained on an enormous unlabelled corpus to
        predict the next token. This is the longest, most
        expensive step. The output is a \emph{base model} that
        is fluent at completing text but does not yet behave as
        an assistant.
  \item \textbf{Instruction tuning.} The base model is shown a
        much smaller corpus of high-quality examples of the kind
        of behaviour it is meant to exhibit --- prompt-and-
        response pairs in which the response is the kind a good
        assistant would give. This step is much shorter and
        cheaper than the pretraining and produces a model that
        responds to instructions rather than continuing the text
        that the instructions appeared in.
  \item \textbf{Preference tuning.} The instruction-tuned model
        is shown pairs of possible responses to the same prompt,
        with a human (or, increasingly, another model) preferring
        one over the other. The model is updated to produce
        responses more like the preferred ones. This is the
        RLHF, DPO, and Constitutional AI work of Chapter~9.
\end{enumerate}

The pretraining is roughly 90\% of the compute and produces
roughly 90\% of the model's capability. The two later stages
adjust the surface behaviour in ways that, at the cost of some
capability, make the model usable. We will look at the later
stages in Chapter~9. The rest of this chapter is about
pretraining itself.

\subsection*{What is in the training data}

The training corpora used for modern Large Language Models are
not published. The general categories are widely understood,
even when the exact composition is not.

Common Crawl, an open project that has been scraping the
public web since 2008, contributes some hundreds of billions of
words once filtered and deduplicated. The model sees ordinary
prose, news articles, blog posts, technical documentation,
forum discussions, and a long tail of less reputable material.
The filtering and deduplication remove the most obviously low-
quality content, but they cannot remove everything, and the
model's behaviour reflects what made it through.

Books contribute another tens to hundreds of billions of words,
depending on how the corpus was assembled. The composition of
the book corpus has been the subject of considerable legal and
ethical controversy --- many of the books in widely-used
training corpora were obtained without the publishers' or
authors' permission, and several lawsuits over this issue are
proceeding through the courts as of this writing.

Code repositories, primarily from GitHub, contribute tens to
hundreds of billions of tokens of source code. This explains
why modern Large Language Models are surprisingly capable
programmers --- they have read more code, in absolute terms,
than any human ever could.

Academic papers, technical documentation, and reference works
contribute a smaller volume but are weighted more heavily in
the training mix because of their per-token information
density. A model that has read the entirety of Wikipedia and
much of arXiv has access to a non-trivial fraction of the
formal knowledge of the species.

The model's capabilities reflect the proportions in this mix.
A model trained on more code is better at code; a model trained
on more academic text is better at academic prose; a model
trained on more conversational data is better at conversation.
Vendors do not, in general, publish the exact proportions, but
the qualitative differences between models in the same
generation are largely traceable to differences in the corpus.

\subsection*{What in-context learning looks like}

The phenomenon Brown and his coauthors named in 2020 is worth
seeing in concrete form, because it is genuinely unusual.

If you take a base model that has been pretrained on next-token
prediction and you give it a prompt that says, in roughly these
words:

\begin{quote}
sea otter $\to$ loutre de mer\\
plush giraffe $\to$ girafe en peluche\\
cheese $\to$
\end{quote}

\noindent the base model will continue the pattern and produce
``fromage''. It will do this even though it was never explicitly
trained on a translation task and even though it has not been
told what the arrows mean. It has, somewhere in its many
billions of parameters and after seeing the pattern in its
training corpus, learned to recognise the format and to produce
a continuation consistent with it. The same trick works for
arithmetic, for word transformation, for question answering, and
for many other tasks if you supply a few examples of the format.

This is the phenomenon that made the field take pretraining
seriously as a path to general capability. A pretrained model
is not a translator, but it is a translator under the right
prompt. It is not a calculator, but it is a calculator under the
right prompt. The same model, with no architectural change and
no further training, performs many tasks. Whatever the model
learned during pretraining, it learned something usefully
broader than ``predict the next word in this corpus''.

We do not, at this writing, have a settled account of why this
works. The leading hypotheses involve some combination of
\emph{induction heads} (specific patterns of attention that the
model develops which copy and continue local patterns),
\emph{implicit Bayesian inference} (the model behaving as if it
is reasoning about which task it is being asked to perform),
and the simple fact that the training corpus contained many
examples of similar pattern-continuation tasks. The honest
position is that we do not know, in any deep sense, why the
trick works as well as it does.

\subsection*{Sampling at the output}

We met the sampling step briefly in Chapter~1. It is worth
revisiting in the context of pretraining because the sampling
choices are part of how the model presents itself to a user.

The pretrained model, given a prompt, produces a probability
distribution over the next token. To produce an actual response,
the model samples from that distribution. The sampling has
several knobs the engineer can turn:

\begin{itemize}
  \item \textbf{Temperature.} A number that, very loosely,
        controls how aggressively the model favours the most
        likely tokens. At temperature zero, the model always
        picks the top of the list (which makes the output
        deterministic but often dull). At temperature one, the
        model samples from the distribution as it stands. Above
        one, the model gives more weight to less likely
        continuations (which makes the output more varied but
        also more error-prone).
  \item \textbf{Top-$k$.} The sampling is restricted to the $k$
        most likely tokens, with everything else zeroed out.
  \item \textbf{Top-$p$ (nucleus).} The sampling is restricted
        to the smallest set of tokens whose combined
        probability exceeds $p$. (This is a more adaptive
        variant of top-$k$.)
\end{itemize}

The default sampling settings of an API have a real effect on
the model's behaviour. A user who cares about reproducibility
should set the temperature to zero or near it. A user who cares
about variety should leave it higher. The same model with
different sampling settings will produce different distributions
of behaviour, and it is reasonable for vendors to tune the
defaults differently for different products.

\section{The engineering consequence}
\addcontentsline{toc}{section}{The engineering consequence}

\paragraph{The training data is the model's worldview.} Almost
every interesting property of a Large Language Model is, in some
sense, a property of its training corpus. A model trained on a
corpus that overrepresents English, the United States, and the
last decade will be best at the kinds of question those subjects
imply. A model trained on a corpus that contained a lot of code
will be better at code. A model trained on a corpus assembled
in 2023 will know nothing of events after 2023. A model trained
on a corpus that included many copies of a particular conspiracy
theory will, when asked, sometimes reproduce it. There is no
fix for this at the model level. The fix is either better data,
later interventions (Chapter~9), or external knowledge sources
(Chapter~10).

\paragraph{The cost of pretraining is structural.} A frontier
pretraining run is a multi-hundred-million-dollar undertaking.
This is not a thing one does lightly. Once a base model is
trained, the economic incentive is to amortise it across as
many products as possible, which is why the same underlying
model often appears under several brand names with different
post-training. A non-technical reader evaluating two LLM
products should ask whether they share a base model; if they
do, the relevant differences are in the post-training and the
deployment, not in the underlying capability.

\paragraph{The base model and the deployed assistant are
different artefacts.} The model OpenAI publishes as ``GPT-4''
in its API is not the GPT-4 base model, but a chain of post-
training steps applied to that base. The assistant's refusals,
its tone, its tendency to add safety disclaimers, its
willingness to perform certain tasks --- all of these are
properties of the post-training, not of the base. A base model,
which OpenAI does not release, would behave differently. This
distinction matters when, for example, a researcher wants to
compare ``what does the model know'' against ``what is the
model willing to say''. The two are, by design, no longer the
same.

\paragraph{Sampling is a product decision.} The defaults for
temperature and other sampling parameters in a vendor's API
shape how the model appears to its users. A vendor who wants
the model to feel ``creative'' will use a higher default
temperature; a vendor who wants the model to feel ``reliable''
will use a lower one. Neither setting is wrong. Both are
choices, and the user of the API can override them. A product
team should think about its sampling defaults the way it would
think about any other user-facing configuration choice.

\section{What goes wrong, and why}
\addcontentsline{toc}{section}{What goes wrong, and why}

\paragraph{Knowledge cut-off.} The model's training data ends
at a particular date. Anything that happened after that date
is, to the model, fiction. A user who asks the model about a
recent event, a recent law, a recently released product, a
recently elected official, will receive an answer that is
plausible but unmoored from the world after the cut-off. The
model usually does not flag this. The fix is either to ask the
model only about pre-cut-off topics or to use a retrieval
system (Chapter~10) that supplies post-cut-off information.

\paragraph{Memorisation and leakage.} The model can, under the
right prompt, reproduce passages from its training corpus
verbatim. For copyrighted material this is a copyright concern.
For personal information this is a privacy concern. For
proprietary data --- in cases where a user has fine-tuned the
model on their own corpus --- this is a confidentiality
concern. There are mitigations (training-time filtering,
output-time filtering, regular auditing), and there is a
research literature that quantifies the risks. The mitigations
are partial, and the risks are not zero.

\paragraph{In-context learning is brittle.} The phenomenon that
gave pretrained models so much of their apparent generality is
also a major source of frustration in production. A prompt that
includes a few examples to specify a format is sensitive to the
choice of examples, the order in which they appear, the
formatting conventions, and the labels. Two prompts that a
human would read as equivalent can produce qualitatively
different outputs. There is a research literature on
\emph{prompt engineering} that catalogues which choices matter,
but no general theory that lets you predict in advance which
prompts will work. The practical advice is to test extensively.

\paragraph{Emergence claims need scrutiny.} Researchers have
reported that certain capabilities --- multi-step arithmetic,
chain-of-thought reasoning, code synthesis --- appear
\emph{suddenly} as the model crosses a threshold of scale,
rather than improving smoothly. The 2022 paper by Jason Wei
and his coauthors, titled \emph{Emergent Abilities of Large
Language Models}, popularised the phrase. A 2023 paper by
Rylan Schaeffer, Brando Miranda, and Sanmi Koyejo, titled
\emph{Are Emergent Abilities of Large Language Models a
Mirage?}, argued that the apparent suddenness is in many cases
a consequence of how the capability is being measured rather
than of how the model is changing. The honest summary is that
some capabilities do appear to scale smoothly when measured
properly, that some still appear to emerge sharply, and that
extrapolating from the current generation to the next on the
basis of either picture is risky.

\section*{Bridges}
\addcontentsline{toc}{section}{Bridges}

This chapter built on Chapter~1 (the cross-entropy objective),
Chapter~2 (the optimisation procedure), Chapter~3 (the deep
neural network), Chapter~4 (the tokenizer and embeddings), and
Chapter~5 (the Transformer architecture). It described what
gets done with all of those components when they are combined
at scale. Everything that follows assumes a pretrained base
model: Chapter~7 discusses what to do when the budget for
pretraining is binding, Chapter~8 discusses how to adapt a
trained base model cheaply, Chapter~9 turns the base model into
an assistant, and Chapter~10 wraps the assistant in retrieval
and tool use. The next chapter --- \emph{Why Bigger Sometimes
Helps} --- is the closer look at the scaling laws and the
recent moves to make them less expensive that we have, so far,
only sketched.
